\documentclass[a4paper, serif, 10pt]{moderncv}

%% moderncv
% style and color
\moderncvstyle{banking}
\moderncvcolor{black}

%% page layout/margins
\usepackage[top=2.5cm, bottom=2.5cm, left=1.5cm, right=1.5cm]{geometry}

\nopagenumbers{}

%% Babel config for Indonesian (Bahasa Indonesia) language
% ref:
% - https://latex3.github.io/babel/guides/locale-indonesian.html
\usepackage[indonesian]{babel}

%% fontspec
\usepackage{fontspec}

\IfFontExistsTF{Linux Libertine}{
  \setmainfont[Ligatures={TeX, Common}, Numbers=OldStyle]{Linux Libertine}
}{}
\IfFontExistsTF{Linux Libertine O}{
  \setmainfont[Ligatures={TeX, Common}, Numbers=OldStyle]{Linux Libertine O}
}{}

%% CTeX
% CJK support, used to show CJK characters in the resume
%
% - fontset=none: disable builtin fontset but instead set the CJK font manually
% - heading=false: disable ctex heading
% - punct=kaiming: use kaiming punctuations styles for CJK
% - scheme=plain: use plain scheme, do not override \normalsize font size
% - space=auto: space settings for CJK characters
%
% ref:
% - http://ctan.mirrorcatalogs.com/language/chinese/ctex/ctex.pdf
\usepackage[UTF8, heading=false, punct=kaiming, scheme=plain, space=auto]{ctex}

\IfFontExistsTF{Noto Serif CJK SC}{
  \setCJKmainfont{Noto Serif CJK SC}
}{}
\IfFontExistsTF{Noto Sans CJK SC}{
  \setCJKsansfont{Noto Sans CJK SC}
}{}

%% line spacing
\usepackage{setspace}
\setstretch{1.125}

%% URL styling - use normal text instead of monospace
\urlstyle{same}

% Auto-underline all links
% Defer until \begin{document} so hyperref is already loaded
\AtBeginDocument{%
  \let\oldhref\href
  \renewcommand{\href}[2]{\oldhref{#1}{\underline{#2}}}%
}

%% Patch \httplink and \httpslink to handle full URLs
% The original moderncv macros always prepend http:// or https:// to URLs.
% This causes issues when the URL already has a protocol (e.g., https://example.com)
% resulting in malformed URLs like https://https://example.com.
% This patch checks if the URL already contains "://" and uses it directly if so.
% Note: We use \str_set:Nx (x-expansion) to expand macros like \@homepage first.
\ExplSyntaxOn
\str_new:N \g_yamlresume_url_str

\RenewDocumentCommand{\httpslink}{O{}m}{
  \str_set:Nx \g_yamlresume_url_str {#2}
  \str_if_in:NnTF \g_yamlresume_url_str {://}
    { \url{#2} }
    { \url{https://#2} }
}

\RenewDocumentCommand{\httplink}{O{}m}{
  \str_set:Nx \g_yamlresume_url_str {#2}
  \str_if_in:NnTF \g_yamlresume_url_str {://}
    { \url{#2} }
    { \url{http://#2} }
}
\ExplSyntaxOff

\name{Rizky Pratama}{}
\title{Data Analyst Berpengalaman dalam SQL, Python, dan Visualisasi Data}
\phone[mobile]{+62 812-3456-7890}
\email{rizky.pratama@email.com}
\homepage{https://rizkypratama.id}

\address{Jl. Senopati No. 45, Jakarta Selatan, DKI Jakarta, Indonesia, 12190}{}{}

\extrainfo{{\small \faLinkedin}\ \href{https://www.linkedin.com/in/rizkypratama}{@rizkypratama} {} {} {} • {} {} {}
{\small \faGithub}\ \href{https://github.com/rizkypratama}{@rizkypratama}}

\begin{document}

\maketitle

\section{Informasi Dasar}

\cvline{}{Data Analyst dengan lebih dari 4 tahun pengalaman mengolah data menjadi insight strategis untuk mendukung pengambilan keputusan bisnis.
%
\begin{itemize}
\item Mahir menggunakan SQL, Python, dan Excel untuk ekstraksi, pembersihan, dan analisis data
\item Berpengalaman membangun dashboard interaktif dengan Tableau dan Power BI untuk memantau performa bisnis secara real-time
\item Mampu menerapkan metode statistik dan A/B testing untuk mengukur dampak inisiatif bisnis
\item Berorientasi pada detail, komunikatif, dan terbiasa bekerja sama dengan tim lintas fungsi
\end{itemize}}
\section{Pendidikan}

\cventry{Agu 2016–Jul 2020}
        {Sarjana, Statistika, Nilai: 3.72}
        {Universitas Indonesia}
        {\url{https://www.ui.ac.id}}
        {}
        {\begin{itemize}
\item Lulus dengan predikat cum laude dan IPK 3.72/4.00
\item Menyelesaikan tugas akhir berjudul "Analisis Faktor yang Mempengaruhi Kepuasan Pelanggan E-commerce Menggunakan Regresi Logistik"
\item Aktif sebagai asisten praktikum mata kuliah Statistika dan Analisis Data
\end{itemize}
\textbf{Mata Kuliah}: Statistika Matematika, Analisis Data Eksploratif, Regresi Linear, Desain Eksperimen, Basis Data dan SQL, Pengantar Machine Learning}
\section{Pengalaman Kerja}

\cventry{Mar 2022–Sekarang}
        {Data Analyst}
        {PT Data Nusantara}
        {\url{https://www.datanusantara.co.id}}
        {}
        {\begin{itemize}
\item Mengembangkan dan memelihara dashboard kinerja perusahaan menggunakan Tableau dan Power BI untuk 5 departemen
\item Menganalisis data penjualan dan perilaku pelanggan dari database SQL dengan volume jutaan baris
\item Menyusun laporan bulanan dan presentasi insight untuk manajemen, mendukung keputusan strategis
\item Berkolaborasi dengan tim produk untuk merancang A/B testing yang meningkatkan konversi sebesar 12\%
\end{itemize}
\textbf{Kata Kunci}: SQL, Python, Tableau, Power BI, A/B Testing}

\cventry{Agu 2020–Feb 2022}
        {Junior Data Analyst}
        {PT Karya Digital Solusi}
        {\url{https://www.karyadigital.co.id}}
        {}
        {\begin{itemize}
\item Membantu proses ETL data penjualan dan stok dari berbagai sumber menggunakan Python dan Google Sheets
\item Membuat visualisasi data untuk kebutuhan laporan mingguan tim marketing dan operasional
\item Melakukan validasi kualitas data serta mendokumentasikan standar pengolahan data
\item Mendukung analisis campaign marketing dan menghitung ROI setiap kanal digital
\end{itemize}
\textbf{Kata Kunci}: Python, ETL, Google Sheets, Visualisasi Data, Marketing Analytics}
\section{Bahasa}

\cvline{Indonesia}{Bahasa Ibu / Bilingual}
\cvline{Inggris}{Tingkat Profesional Penuh \hfill \textbf{Kata Kunci}: TOEFL ITP 600, IELTS 7.0}
\section{Keahlian}

\cvline{Analisis Data}{Ahli \hfill \textbf{Kata Kunci}: SQL, Python, Pandas, NumPy, Excel, Google Sheets}
\cvline{Statistika}{Mahir \hfill \textbf{Kata Kunci}: Regresi, Uji Hipotesis, A/B Testing, Forecasting, R}
\cvline{Visualisasi Data}{Mahir \hfill \textbf{Kata Kunci}: Tableau, Power BI, Looker Studio, Matplotlib, Seaborn}
\section{Penghargaan}

\cventry{Des 2023}
        {PT Data Nusantara}
        {Data Analyst Terbaik Tahun 2023}
        {}
        {}
        {Penghargaan internal atas kontribusi dalam membangun sistem pelaporan data yang menghemat waktu analisis hingga 30\%.}
\section{Sertifikat}

\cventry{Mar 2021}
        {Google}
        {Google Data Analytics Professional Certificate}
        {\url{https://www.coursera.org/professional-certificates/google-data-analytics}}
        {}
        {}

\cventry{Jun 2022}
        {HackerRank}
        {SQL Advanced Certificate}
        {\url{https://www.hackerrank.com/certificates/sql-advanced}}
        {}
        {}
\section{Proyek}

\cventry{Mar 2023–Jun 2023}
        {Dashboard interaktif untuk menganalisis tren penjualan dan perilaku pelanggan.}
        {Dashboard Penjualan E-commerce}
        {\url{https://public.tableau.com/app/profile/rizky.pratama}}
        {}
        {\begin{itemize}
\item Dirancang menggunakan Tableau dengan koneksi langsung ke Google BigQuery
\item Menampilkan metrik utama seperti GMV, conversion rate, dan retention rate secara real-time
\item Memungkinkan manajemen memantau performa harian per kategori produk dan wilayah
\end{itemize}
\textbf{Kata Kunci}: Tableau, BigQuery, E-commerce, KPI}

\cventry{Sep 2022–Des 2022}
        {Proyek analisis sentimen ulasan pelanggan menggunakan Python dan machine learning.}
        {Analisis Sentimen Ulasan Pelanggan}
        {\url{https://github.com/rizkypratama/sentiment-analysis}}
        {}
        {\begin{itemize}
\item Melakukan scraping data ulasan dari platform e-commerce dan membersihkan teks dengan Python
\item Membangun model klasifikasi sentimen menggunakan scikit-learn dengan akurasi 87\%
\item Menyajikan insight utama mengenai faktor kepuasan dan keluhan pelanggan
\end{itemize}
\textbf{Kata Kunci}: Python, Machine Learning, NLP, scikit-learn}
\section{Minat}

\cvline{Data Science}{Machine Learning, Statistika, Data Storytelling}
\cvline{Olahraga}{Lari, Badminton, Hiking}
\section{Kegiatan Sukarela}

\cventry{Jan 2021–Des 2021}
        {Relawan Analisis Data}
        {Yayasan Pustaka Cerdas}
        {\url{https://pustakacerdas.or.id}}
        {}
        {\begin{itemize}
\item Membantu yayasan mengelola data donatur dan penerima bantuan menggunakan spreadsheet
\item Menyusun laporan tahunan dengan visualisasi yang mudah dipahami oleh pengurus yayasan
\end{itemize}}

\end{document}